% APÊNDICES -------------------------------------------------------------------

\begin{apendicesenv}

\chapter{Códigos CNAE e CBO que definem o recorte de tecnologia}
\label{apend:cbos}

Este apêndice documenta o recorte aplicado, discutido na \autoref{sec:recorte}. Um vínculo entra na pesquisa se o estabelecimento tem CNAE de atividade de TI \textbf{ou} se a ocupação tem CBO de TI. As duas condições são mantidas como rótulos independentes em cada registro da camada tratada, o que permite distinguir as duas lentes nas análises.

\section{Atividades econômicas (CNAE 2.0)}

A \autoref{quadro:cnae} relaciona as subclasses CNAE que caracterizam a empresa como sendo do setor de TI.

\begin{quadro}[H]
    \centering
    \caption{Subclasses CNAE que definem o setor de tecnologia}
    \label{quadro:cnae}
    \footnotesize
    \renewcommand{\arraystretch}{1.25}
    \begin{tabular}{lp{10.0cm}}
        \toprule
        \rowcolor{black!8}
        \textbf{CNAE} & \textbf{Atividade} \\
        \midrule
        6201500 & Desenvolvimento de programas de computador sob encomenda \\
        6201501 & Desenvolvimento de programas de computador sob encomenda \\
        6201502 & \textit{Web design} \\
        6202300 & Desenvolvimento e licenciamento de programas customizáveis \\
        6203100 & Desenvolvimento e licenciamento de programas não customizáveis \\
        6204000 & Consultoria em tecnologia da informação \\
        6209100 & Suporte técnico, manutenção e outros serviços em TI \\
        6311900 & Tratamento de dados, hospedagem na internet e atividades relacionadas \\
        6319400 & Portais, provedores de conteúdo e outros serviços de informação \\
        6399200 & Outras atividades de prestação de serviços de informação \\
        \bottomrule
    \end{tabular}
    \\[4pt]
    Fonte: Autoria própria, a partir de \citeonline{CNAE2024}.
\end{quadro}

\section{Famílias ocupacionais (CBO 2002)}

A seleção de ocupações é feita por \textbf{família} --- os quatro primeiros dígitos do código ---, mais três códigos avulsos, e menos cinco códigos excluídos. A \autoref{quadro:familias} apresenta essa definição com a justificativa de cada decisão.

\begin{quadro}[H]
    \centering
    \caption{Famílias e códigos CBO que definem a ocupação de tecnologia}
    \label{quadro:familias}
    \footnotesize
    \renewcommand{\arraystretch}{1.25}
    \begin{tabular}{llp{7.6cm}}
        \toprule
        \rowcolor{black!8}
        \textbf{Código} & \textbf{Situação} & \textbf{Família ocupacional e justificativa} \\
        \midrule
        1236   & Incluída & Gerentes de pesquisa e desenvolvimento: direção de serviços de informática. \\
        1425   & Incluída & Gerentes de tecnologia da informação. \\
        2031   & Incluída & Pesquisadores das ciências da computação; demais pesquisadores da família foram excluídos. \\
        2112   & \textbf{Incluída na revisão} & Estatísticos: família sob a qual analistas e cientistas de dados são registrados. Ver \autoref{subsec:revisao-recorte}. \\
        2122   & Incluída & Engenheiros em computação. \\
        2123   & Incluída & Administradores de redes, sistemas e bancos de dados. \\
        2124   & Incluída & Analistas de tecnologia da informação: núcleo do setor. \\
        3171   & Incluída & Técnicos de desenvolvimento de sistemas e aplicações. \\
        3172   & Incluída & Técnicos em operação e monitoração de computadores. \\
        \midrule
        313220 & Avulso   & Técnico em manutenção de equipamentos de informática. \\
        313305 & Avulso   & Técnico em telecomunicações com aderência a redes de dados. \\
        142135 & Avulso   & Gerente de rede e de sistemas (código recente). \\
        \midrule
        203110 & \textbf{Excluído} & Pesquisador em ciências da terra e meio ambiente. \\
        203115 & \textbf{Excluído} & Pesquisador em física. \\
        203120 & \textbf{Excluído} & Pesquisador em matemática. \\
        203125 & \textbf{Excluído} & Pesquisador em química. \\
        317115 & \textbf{Excluído} & Programador de máquinas-ferramenta com comando numérico: ocupação industrial. \\
        \bottomrule
    \end{tabular}
    \\[4pt]
    Nota: a exclusão incide apenas sobre a lente de ocupação. Um pesquisador em física
    empregado em empresa de TI permanece no recorte pela lente de setor. \\
    Fonte: Autoria própria, a partir de \citeonline{CBO2024}.
\end{quadro}

\section{Códigos observados no estoque}

O \autoref{quadro:cbos} relaciona todos os códigos de seis dígitos efetivamente encontrados no estoque do último ano disponível, com a descrição oficial, a área atribuída e o número de vínculos. O quadro é \textbf{gerado a partir da camada Gold}, e não digitado: qualquer revisão do recorte se reflete nele automaticamente. É assim que se evita o problema descrito na \autoref{sec:achados-qualidade}, em que uma lista mantida à mão passou a divergir do recorte efetivamente aplicado.

% Gerado por scripts/gerar_tabelas.py --- não editar à mão.
\begin{quadro}[htbp]
    \centering
    \caption{Códigos CBO do recorte de tecnologia observados no estoque de 2025}
    \label{quadro:cbos}
    \footnotesize
    \renewcommand{\arraystretch}{1.25}
    \begin{tabular}{llp{6.0cm}lr}
        \toprule
        \rowcolor{black!8}
        \textbf{Família} & \textbf{CBO} & \textbf{Ocupação} & \textbf{Área} & \textbf{Vínculos 2025} \\
        \midrule
        \texttt{1236} & \texttt{123605} & Diretor de Servicos de Informatica & Gestão e Direção & 4.649 \\
        \texttt{1421} & \texttt{142135} & Código recente, sem descrição no dicionário do MTE & Gestão e Direção & 432 \\
        \texttt{1425} & \texttt{142505} & Gerente de Rede & Gestão e Direção & 8.859 \\
        \texttt{1425} & \texttt{142510} & Gerente de Desenvolvimento de Sistemas & Gestão e Direção & 20.568 \\
        \texttt{1425} & \texttt{142515} & Gerente de Producao de Tecnologia da Informacao & Gestão e Direção & 8.526 \\
        \texttt{1425} & \texttt{142520} & Gerente de Projetos de Tecnologia da Informacao & Gestão e Direção & 29.512 \\
        \texttt{1425} & \texttt{142525} & Gerente de Seguranca de Tecnologia da Informacao & Gestão e Direção & 2.274 \\
        \texttt{1425} & \texttt{142530} & Gerente de Suporte Tecnico de Tecnologia da Informac & Gestão e Direção & 12.027 \\
        \texttt{1425} & \texttt{142535} & Tecnólogo em Gestão da Tecnologia da Informação & Gestão e Direção & 5.960 \\
        \texttt{2031} & \texttt{203105} & Pesquisador em Ciencias da Computacao e Informatica & Engenharia e Pesquisa & 3.403 \\
        \texttt{2112} & \texttt{211205} & Estatistico & Estatística e Ciência de Dados & 4.528 \\
        \texttt{2112} & \texttt{211210} & Estatistico (Estatistica Aplicada) & Estatística e Ciência de Dados & 388 \\
        \texttt{2112} & \texttt{211215} & Estatistico Teorico & Estatística e Ciência de Dados & 123 \\
        \texttt{2112} & \texttt{211220} & Código recente, sem descrição no dicionário do MTE & Estatística e Ciência de Dados & 2.860 \\
        \texttt{2122} & \texttt{212205} & Engenheiro de Aplicativos em Computacao & Engenharia e Pesquisa & 12.059 \\
        \texttt{2122} & \texttt{212210} & Engenheiro de Equipamentos em Computacao & Engenharia e Pesquisa & 838 \\
        \texttt{2122} & \texttt{212215} & Engenheiros de Sistemas Operacionais em Computacao & Engenharia e Pesquisa & 12.601 \\
        \texttt{2123} & \texttt{212305} & Administrador de Banco de Dados & Infraestrutura e Dados & 14.060 \\
        \texttt{2123} & \texttt{212310} & Administrador de Redes & Infraestrutura e Dados & 7.043 \\
        \texttt{2123} & \texttt{212315} & Administrador de Sistemas Operacionais & Infraestrutura e Dados & 11.529 \\
        \texttt{2123} & \texttt{212320} & Administrador em Segurança da Informação & Infraestrutura e Dados & 15.129 \\
        \texttt{2124} & \texttt{212405} & Analista de Desenvolvimento de Sistemas & Desenvolvimento & 294.260 \\
        \texttt{2124} & \texttt{212410} & Analista de Redes e de Comunicacao de Dados & Desenvolvimento & 53.475 \\
        \texttt{2124} & \texttt{212415} & Analista de Sistemas de Automacao & Desenvolvimento & 11.587 \\
        \texttt{2124} & \texttt{212420} & Analista de Suporte Computacional & Desenvolvimento & 110.598 \\
        \texttt{2124} & \texttt{212425} & Código recente, sem descrição no dicionário do MTE & Desenvolvimento & 6.965 \\
        \texttt{2124} & \texttt{212430} & Código recente, sem descrição no dicionário do MTE & Desenvolvimento & 13.651 \\
        \texttt{3132} & \texttt{313220} & Tecnico em Manutencao de Equipamentos de Informatica & Suporte e Operação & 48.372 \\
        \texttt{3133} & \texttt{313305} & Tecnico de Comunicacao de Dados & Infraestrutura e Dados & 7.049 \\
        \texttt{3171} & \texttt{317105} & Programador de Internet & Desenvolvimento & 5.213 \\
        \texttt{3171} & \texttt{317110} & Programador de Sistemas de Informacao & Desenvolvimento & 90.858 \\
        \texttt{3171} & \texttt{317120} & Programador de Multimidia & Desenvolvimento & 2.115 \\
        \texttt{3172} & \texttt{317205} & Operador de Computador (Inclusive Microcomputador) & Suporte e Operação & 31.761 \\
        \texttt{3172} & \texttt{317210} & Tecnico de Apoio ao Usuario de Informatica (Helpdesk & Suporte e Operação & 90.619 \\
        \bottomrule
    \end{tabular}
    \\[4pt]
    Ocupações sem descrição são códigos em uso que não constam dos dicionários oficiais do MTE. Fonte: Autoria própria, a partir da camada Gold.
\end{quadro}


\chapter{Consultas de verificação}
\label{apend:consultas}

As consultas a seguir podem ser executadas por qualquer pessoa, sem acesso à infraestrutura do projeto, bastando um cliente DuckDB com a extensão \texttt{httpfs}. Elas leem os dados publicados diretamente por HTTP.

\section{Reprodutibilidade: o dado bruto reproduz a camada publicada?}

Esta consulta aplica o recorte de tecnologia ao microdado bruto de um mês e compara o resultado com a camada tratada publicada. Diferença zero significa que o recorte é reproduzível a partir da fonte, conforme discutido na \autoref{sec:reprodutibilidade}.

\begin{lstlisting}[language=SQL, basicstyle=\ttfamily\scriptsize, breaklines=true,
                   frame=single, captionpos=b,
                   caption={Verificação de reprodutibilidade do recorte},
                   label={lst:reprodutibilidade}]
WITH rotulado AS (
    SELECT subclasse,
           lpad(regexp_replace(trim(cbo2002ocupacao), '[^0-9]', '', 'g'), 6, '0') AS cbo6
    FROM read_parquet('hf://datasets/Gianpedro/bronze_caged/caged_mov/ano=2025/mes=6/*.parquet')
),
cru AS (
    SELECT count(*) AS linhas FROM rotulado
    WHERE ( ( substr(cbo6, 1, 4) IN ('1236','1425','2031','2112','2122','2123','2124','3171','3172')
              OR cbo6 IN ('313220','313305','142135') )
            AND cbo6 NOT IN ('203110','203115','203120','203125','317115') )
       OR subclasse IN ('6201500','6201501','6201502','6202300','6203100',
                        '6204000','6209100','6311900','6319400','6399200')
),
publicado AS (
    SELECT count(*) AS linhas
    FROM read_parquet('hf://datasets/Gianpedro/caged-tecnologia/caged_mov/ano_particao=2025/mes_particao=6/*.parquet')
)
SELECT cru.linhas AS no_dado_cru, publicado.linhas AS na_camada_publicada,
       cru.linhas - publicado.linhas AS diferenca
FROM cru, publicado;
\end{lstlisting}

\section{Procedência de uma tradução}

Esta consulta recupera a origem de qualquer descrição aplicada pelo \textit{pipeline}: de qual planilha de \textit{layout}, de qual aba, de qual caminho no servidor público e em que data a extração foi feita.

\begin{lstlisting}[language=SQL, basicstyle=\ttfamily\scriptsize, breaklines=true,
                   frame=single, captionpos=b,
                   caption={Procedência das traduções de uma coluna codificada},
                   label={lst:procedencia}]
SELECT DISTINCT tabela, coluna, planilha, aba, caminho_ftp, extraido_em
FROM read_parquet('hf://datasets/Gianpedro/bronze_caged/dicionarios.parquet')
WHERE coluna = 'sexo';
\end{lstlisting}

\end{apendicesenv}
